\chapter{Related Work and Conclusion}
\label{sec:related}

Sulzmann and Lu observed that many POSIX regular expression
implementations were incorrect. To address this problem, they proposed
a method based on Brzozowski's derivatives~\citep{brzozowski1964},
which constructs a full parse tree according to the longest-leftmost
disambiguation policy. They also provided a proof of correctness
for their approach.
In separate work, they investigated submatching, an area in which
existing approaches predominantly rely on greedy disambiguation.
Greedy matching is closely tied to the structure of the regular
expression and can be implemented efficiently by translating the
expression into a nondeterministic finite automaton (NFA). 
Examples include the general parsing approach of
\citet{frischcardelli2004} and Laurikari's tagged NFAs \citet{laurikari2000} for
submatching. Applying these NFA-based approaches to POSIX semantics,
however, requires additional bookkeeping to select the
longest-leftmost result, as observed by
\citet{sulzmannlu2012}. Instead, their approach to submatching
uses Antimirov's partial derivatives directly, resulting in a
non-POSIX counterpart to their full-parse-tree construction
\citep{sulzmannlu2012,antimirov1996}.

While the original papers provided Haskell implementations, their
practical realisation in a systems programming language had not
been investigated. This thesis addressed this gap by implementing
all four parser variants in Rust and empirically evaluating their
performance and memory consumption on the S2ELab
corpus~\citep{s2elab-regexeval-dataset} and a synthetic stress-test
dataset.

Across the datasets, the implementations had the following
performance ordering, from the slowest to the fastest parser:
\texttt{deriv\_std\_rec}, \texttt{deriv\_std\_loop},
\texttt{deriv\_bc}, \texttt{pderiv\_std}, and
\texttt{pderiv\_bc}. In particular, \texttt{deriv\_bc} was four
to five times faster than the standard derivative parser while
using roughly two orders of magnitude less peak memory.
Partial derivatives proved more efficient for submatching on
realistic inputs, whereas bit-coding substantially benefited the
full-derivative parser but made little difference for partial
derivatives. The relative performance nevertheless depends on
the input, particularly on the growth of the residual set, so
neither family was shown to be uniformly superior.

Beyond derivative-based approaches, \citet{laurikari2000} attach
submatch information to automaton transitions, as in Kuklewicz's
\texttt{regex-tdfa}~\citep{kuklewicz_tdfa}. It is one of three tools
\citet{sulzmannlu2014} benchmark against in their own paper, not with
the general construction presented there but with a second, adapted
Haskell implementation that records only the last match of each
Kleene star. The bit-coded parsers here implement that general
construction directly, without this adaptation.
Production engines such as
RE2~\citep{re2} and Rust's \texttt{regex}~\citep{rust_regex} compile
to automata for linear-time matching but do not build the full parse
trees this thesis targets.

For future work, three directions follow from these results. 
First, \texttt{deriv\_bc}
still trails the partial-derivative variants on several
large-scale patterns, partly due to repeatedly simplifying a
shared expression against itself rather than stepping independent
residuals. Sulzmann and Lu report memory problems in their
bit-coded implementation and propose computing the POSIX path
separately before constructing the parse tree with a
space-efficient tree builder~\citep{sulzmannlu2014}. A related
improvement for \texttt{deriv\_bc} would be to defer bit construction
until decoding. 
Second, the supported language could be extended to include
operations currently rejected by the frontend, such as anchors,
bounded lookaround, and backreferences, as discussed in
Chapter~\ref{sec:frontend}. Lu's partial-derivative
implementation~\citep{lu_regex_pderiv} provides a useful reference
for such extensions.
Additionally, the current trace instrumentation could be developed
into a proper diagnostic tool identifying which subexpressions are
responsible for a match or timeout, helping to localise remaining
performance problems or providing verbose information on where and
why a given input failed to match.